\section{Schema Boundary}

RPLAN stores plan identity and assignments. RCTX stores graph, population,
canonical unit order, and context needed for verification. Certificates and
manifests record declared profiles, verifier results, hashes, lineage, and
algorithm metadata.

\begin{table}[htbp]
\centering
\small
\begin{tabular}{p{0.20\linewidth}p{0.35\linewidth}p{0.31\linewidth}}
\toprule
Artifact & Owns & Must not own \\
\midrule
RPLAN & final assignments and plan identity & graph/population context \\
RCTX & units, graph, population, canonical order & algorithm-specific proof \\
Certificate & profile checks, verifier results, lineage & legal compliance claim \\
Manifest & files, hashes, package inventory & substantive validity by itself \\
\bottomrule
\end{tabular}
\caption{RPLAN/RCTX fixed-point artifact responsibilities.}
\end{table}

The fixed point is a shared final-plan contract, not a new algorithmic claim.
Construction, exact optimization, local search, and Pareto comparison may carry
different upstream evidence, but their final packages must agree on the same
artifact vocabulary.

The atlas framing requires both positive and negative verifier paths to be
concrete. A passing package shows the plan, context, certificate, manifest, and
hash chain agreeing. A failing package should identify the exact broken
invariant: missing context, order mismatch, stale hash, unsupported profile, or
broken lineage pointer. That failure reason is part of the audit product rather
than an implementation footnote.

\begin{table}[htbp]
\centering
\small
\begin{tabular}{p{0.27\linewidth}p{0.38\linewidth}p{0.22\linewidth}}
\toprule
Invariant & Required check & Failure mode \\
\midrule
Canonical unit order & RPLAN assignments align with RCTX units & order mismatch \\
Hash integrity & manifest hashes match files & tamper or stale file \\
Profile scope & certificate names declared checks & unsupported claim \\
Lineage & algorithm/run ids link to parent artifacts & broken provenance \\
\bottomrule
\end{tabular}
\caption{Fixed-point invariants for verifier-facing packages.}
\end{table}

\begin{table}[htbp]
\centering
\small
\begin{tabular}{p{0.22\linewidth}p{0.28\linewidth}p{0.34\linewidth}}
\toprule
Family & Example package & Evidence retained before U.20 \\
\midrule
Construction & T.14--T.17 golden packages & construction summaries, witnesses,
and final assignments \\
Exact search & U.16 and U.17 packages & solve reports, separation/column
events, and model-relative certificates \\
Local search & U.18 package & parent plan, candidate move, child plan, and
improvement status \\
Pareto search & U.19 selected-frontier package & frontier row, selected index,
objectives, and export metadata \\
Audit benchmark & U.20 grid packages & manifest, certificate, context, hashes,
and negative fixture catalog \\
\bottomrule
\end{tabular}
\caption{Public package families currently exercising the fixed point.}
\end{table}

The common package shape is deliberately small. A typical public package has a
plan file, a context file, an audit-certificate JSON file, and a manifest JSON
file. Method-specific files may sit beside those artifacts, but the
verifier-facing contract is the same: the manifest names the files and hashes,
the certificate binds the plan and context under a declared profile, and the
lineage record says which algorithm run produced the exported plan.
